\section{Related work}
\label{sec:related-work}

\paragraph{Certified bounds for permutation puzzles.}
Group actions, Cayley graphs, and subgroup methods provide the standard
language for permutation puzzles~\cite{joyner2008}.  Pattern databases and
memory-efficient heuristic search made large puzzle spaces computationally
accessible~\cite{korf1997,culberson1998}.  Kunkle and Cooperman used
computational group methods to establish a 26-move upper bound for the Rubik's
Cube~\cite{kunkle2007}; Rokicki, Kociemba, Davidson, and Dethridge later
certified that the diameter in the half-turn metric is 20~\cite{rokicki2013}.
These works exemplify the separation between a large discovery computation
and a smaller, auditable verification argument.  The present computation uses
the same high-level discipline, but its finite obligation is a pair of
maximal-layer products in a Megaminx subgroup chain.

\paragraph{Learned heuristic search.}
DeepCubeA combines a learned cost-to-go function with batched heuristic search
and solves random instances of several puzzles, including the Rubik's
Cube~\cite{agostinelli2019}.  EfficientCube demonstrates that comparatively
simple self-supervision from goal-rooted scrambles can support neural-guided
beam search~\cite{takano2023}.  More recent work uses neural distance models
and massively parallel beam search on larger Rubik's Cubes
\cite{chervov2025}.  Our training signal and model architecture are secondary
to the certificate result: the models need only rank sufficiently many promising actions
to find a valid word for every member of a fixed exhaustive frontier.

\paragraph{Beam portfolios and batched learned heuristics.}
Restarting a learned beam search with progressively larger widths has appeared
in combinatorial optimization~\cite{cohenbeck2021}.  GPU batching of expensive
learned heuristics is likewise established in heuristic-search research
\cite{greco2022,futuhisturtevant2026}.  Accordingly, we describe our procedure
as a \emph{beam cascade} or staged search portfolio, not as a new beam-search
algorithm.  Its role in our computation is proof-oriented budget allocation: each
stage searches only unresolved members of an already exhaustive family, and
each success is persisted as a model-independent certificate.

\paragraph{Megaminx computations.}
The coordinates, phase numbering, move conventions, and reported global bound
used here originate in the software and accompanying analysis of
Botz~\cite{botz2026dataset,botzsource}.  The public history begins with the
116-move chain~\cite{botz2025bound116}.  A later table correction confirmed,
rather than improved, that value~\cite{botz2025correction}; two subsequent
reports reduced the last-phase contribution $26\to25\to24$, giving
$116\to115\to114$~\cite{rokicki2025bound115,rokicki2025bound114,botz2026bound114}.
These are external software and forum reports, not computations regenerated in
this paper.  The separate phase-1 result proves that allowing
the first block target to vary over all $20$ corners lowers its worst-case cost
from 10 to 9~\cite{kuznetsov2026phase1}.  Neither that result nor the present
paper reruns all external computations behind the reported 114-move chain;
the local theorems are unconditional, whereas the global arithmetic is not.
